\documentclass{article}
\usepackage[margin=1in, a4paper]{geometry}
\usepackage{amsmath, amssymb, amsfonts}
\usepackage{graphicx}
\usepackage{xcolor}
\usepackage{listings}
\usepackage{booktabs}
\usepackage[utf8]{inputenc}
\usepackage[T1]{fontenc}
\usepackage{hyperref}
\hypersetup{colorlinks=true, urlcolor=blue, linkcolor=blue}
\usepackage{enumitem}
\usepackage{float}

\definecolor{codegreen}{rgb}{0,0.6,0}
\definecolor{codegray}{rgb}{0.5,0.5,0.5}
\definecolor{codepurple}{rgb}{0.58,0,0.82}
\definecolor{backcolour}{rgb}{0.99,0.99,0.99}
% Professional color palette for academic publishing
\definecolor{myblue}{cmyk}{0.8,0.4,0,0.1}
\definecolor{mygreen}{cmyk}{0.7,0,0.5,0.2}
\definecolor{myorange}{cmyk}{0,0.5,0.8,0.1}
\definecolor{mygray}{cmyk}{0,0,0,0.4}
\definecolor{arrowcolor}{cmyk}{0.9,0.5,0,0.3}
\definecolor{nodecolor}{cmyk}{0.6,0.2,0,0.05}
\definecolor{framecolor}{cmyk}{0.4,0.1,0,0.1}

\lstdefinestyle{mystyle}{
    backgroundcolor=\color{backcolour},   
    commentstyle=\color{codegreen},
    keywordstyle=\color{magenta},
    numberstyle=\tiny\color{codegray},
    stringstyle=\color{codepurple},
    basicstyle=\ttfamily\footnotesize,
    breakatwhitespace=false,         
    breaklines=true,                 
    captionpos=b,                    
    keepspaces=true,                 
    numbers=left,                    
    numbersep=5pt,                  
    showspaces=false,                
    showstringspaces=false,
    showtabs=false,                  
    tabsize=2
}
\lstset{style=mystyle}

\title{The Single vs. Dual Sourcing Problem}
\author{The Logistics \& Supply Chain Problem-Solving Playbook}
\date{}

\begin{document}
\maketitle

\section{The Single vs. Dual Sourcing Problem}

\subsection{The Scenario: Balancing Risk and Efficiency in Strategic Procurement}

Consider MediTech Manufacturing, a company producing critical medical devices that requires specialized electronic components for their life-support equipment. The procurement team faces a fundamental strategic decision: should they source these critical components from a single trusted supplier (single sourcing) or split their requirements between two suppliers (dual sourcing)?

With single sourcing, MediTech can leverage economies of scale, negotiate better pricing through volume commitments, and build deep collaborative relationships with their chosen supplier. However, this strategy exposes them to significant supply chain risks---if their sole supplier experiences disruptions, faces quality issues, or encounters financial difficulties, MediTech's entire production line could halt.

Dual sourcing offers a different value proposition: by splitting orders between two suppliers, MediTech reduces dependency risk and creates competitive pressure that may drive down costs. However, this approach requires more complex supplier management, potentially sacrifices volume discounts, and demands greater administrative overhead.

This problem becomes even more complex when considering multiple components with varying criticality levels, different supplier capabilities, seasonal demand fluctuations, and the need to maintain consistent quality standards across all sourcing decisions.

\subsection{Tier 1: The Pen \& Paper Method (Multi-Objective Mathematical Formulation)}

The single vs. dual sourcing decision can be formulated as a multi-objective optimization problem that simultaneously minimizes cost and risk while ensuring supply reliability.

\textbf{Mathematical Model}

Let us define the following sets and parameters:

\textbf{Sets:}
\begin{align}
I &= \{1, 2, \ldots, n\} \quad \text{set of components} \\
J &= \{1, 2, \ldots, m\} \quad \text{set of potential suppliers} \\
T &= \{1, 2, \ldots, T\} \quad \text{set of time periods}
\end{align}

\textbf{Parameters:}
\begin{align}
d_{it} &= \text{demand for component } i \text{ in period } t \\
c_{ij} &= \text{unit cost of component } i \text{ from supplier } j \\
r_{j} &= \text{reliability score of supplier } j \text{ (0-1 scale)} \\
q_{ij} &= \text{quality score of component } i \text{ from supplier } j \\
v_{ij} &= \text{volume discount factor for component } i \text{ from supplier } j \\
M &= \text{large positive constant}
\end{align}

\textbf{Decision Variables:}
\begin{align}
x_{ijt} &\quad \text{quantity of component } i \text{ sourced from supplier } j \text{ in period } t \\
y_{ij} &\quad \text{1 if supplier } j \text{ is selected for component } i, \text{0 otherwise} \\
z_{i} &\quad \text{1 if component } i \text{ uses dual sourcing, 0 if single sourcing}
\end{align}

\textbf{Objective Functions:}

\textbf{Minimize Total Cost:}
\begin{align}
f_1 = \sum_{i \in I} \sum_{j \in J} \sum_{t \in T} c_{ij} \cdot v_{ij} \cdot x_{ijt}
\end{align}

\textbf{Maximize Supply Chain Reliability:}
\begin{align}
f_2 = \sum_{i \in I} \sum_{j \in J} r_j \cdot y_{ij} \cdot \left(\sum_{t \in T} x_{ijt}\right)
\end{align}

\textbf{Constraints:}

\textbf{Demand Satisfaction:}
\begin{align}
\sum_{j \in J} x_{ijt} = d_{it} \quad \forall i \in I, t \in T
\end{align}

\textbf{Supplier Selection Logic:}
\begin{align}
x_{ijt} \leq M \cdot y_{ij} \quad \forall i \in I, j \in J, t \in T
\end{align}

\textbf{Sourcing Strategy Constraint:}
\begin{align}
\sum_{j \in J} y_{ij} = 1 + z_i \quad \forall i \in I
\end{align}

\textbf{Quality Requirements:}
\begin{align}
\sum_{j \in J} q_{ij} \cdot y_{ij} \geq Q_{\min} \quad \forall i \in I
\end{align}

\begin{figure}[htbp]
\centering
\includegraphics[width=\textwidth]{../figures-png/line89.fig1.png}
\caption{Enhanced Multi-Sourcing Procurement Problem with Node-Based Decision Structure and Professional Styling}
\end{figure}

\textbf{Concrete Example}

Consider a simplified scenario with 2 components and 3 potential suppliers over 4 quarters:

\begin{itemize}
\item Component 1 (Critical): Demand = [100, 120, 110, 130] units per quarter
\item Component 2 (Standard): Demand = [200, 180, 220, 210] units per quarter
\end{itemize}

Supplier data:
\begin{center}
\begin{tabular}{|c|c|c|c|}
\hline
\textbf{Supplier} & \textbf{Cost (Comp 1)} & \textbf{Cost (Comp 2)} & \textbf{Reliability} \\
\hline
A & \$15 & \$8 & 0.85 \\
B & \$18 & \$7 & 0.92 \\
C & \$16 & \$9 & 0.88 \\
\hline
\end{tabular}
\end{center}

Solving the multi-objective formulation yields:
\begin{itemize}
\item Component 1: Dual sourcing with 60\% from Supplier B, 40\% from Supplier C
\item Component 2: Single sourcing from Supplier B
\item Total annual cost: \$8,760 (vs. \$8,200 for pure single sourcing)
\item Reliability improvement: 15\% reduction in supply disruption risk
\end{itemize}

\subsection{Tier 2: The Classic Heuristic (Priority-Based Allocation Algorithm)}

A priority-based heuristic approach evaluates sourcing decisions by assigning weighted scores to multiple criteria and systematically allocating components to sourcing strategies.

\textbf{Algorithm Explanation}

The Priority-Based Sourcing Algorithm uses a composite scoring function that weighs cost efficiency, supply risk, and operational complexity. Components are ranked by criticality, and sourcing decisions are made sequentially based on threshold values.

Time Complexity: O(n $\times$ m $\times$ log m) where n is the number of components and m is the number of suppliers.

\textbf{Pseudocode}

\lstinputlisting[language=Python, caption=Priority-Based Sourcing Algorithm]{.././codes-python/line89.py1.py}


\begin{figure}[htbp]
\centering
\includegraphics[width=\textwidth]{../figures-png/line89.fig2.png}
\caption{Priority-Based Sourcing Algorithm Architecture}
\end{figure}

\textbf{Concrete Example}

Using the same data from Tier 1, the priority-based algorithm processes components as follows:

Input data:
\begin{itemize}
\item Component 1: Criticality = 0.9, Annual Demand = 460 units
\item Component 2: Criticality = 0.6, Annual Demand = 810 units
\end{itemize}

Composite scores calculation:
\begin{center}
\begin{tabular}{|c|c|c|c|}
\hline
\textbf{Supplier} & \textbf{Component 1 Score} & \textbf{Component 2 Score} & \textbf{Ranking} \\
\hline
A & 0.72 & 0.78 & 2nd for Comp1, 1st for Comp2 \\
B & 0.81 & 0.75 & 1st for Comp1, 2nd for Comp2 \\
C & 0.74 & 0.71 & 3rd for Comp1, 3rd for Comp2 \\
\hline
\end{tabular}
\end{center}

Algorithm output:
\begin{itemize}
\item Component 1: Dual sourcing (70\% Supplier B, 30\% Supplier A)
\item Component 2: Single sourcing (100\% Supplier A)
\item Processing time: 0.02 seconds for 2 components, 3 suppliers
\item Decision confidence: 87\% based on score separations
\end{itemize}

\subsection{Tier 3: The Advanced Algorithm (Ant Colony Optimization)}

Ant Colony Optimization (ACO) models the sourcing problem as a graph traversal challenge where artificial ants explore different sourcing combinations, leaving pheromone trails that guide future search towards optimal solutions.

\textbf{Theory Explanation}

In the sourcing context, each ant represents a complete sourcing solution. Ants traverse a construction graph where nodes represent supplier-component assignments and edges represent transitions between assignments. The pheromone intensity on edges reflects the quality of past solutions, while heuristic information guides ants toward promising regions of the solution space.

The algorithm balances exploration (trying new sourcing combinations) with exploitation (leveraging proven successful patterns) through pheromone updating and evaporation mechanisms.

\textbf{Pseudocode}

\lstinputlisting[language=Python, caption=Ant Colony Optimization for Sourcing]{.././codes-python/line89.py2.py}

\begin{figure}[htbp]
\centering
\includegraphics[width=\textwidth]{../figures-png/line89.fig3.png}
\caption{Ant Colony Optimization Architecture for Sourcing Decisions}
\end{figure}

\textbf{Concrete Example}

Applied to our medical device manufacturing scenario with 5 components and 4 suppliers over 100 iterations:

ACO Parameters:
\begin{itemize}
\item Number of ants: 30
\item $\alpha$ (pheromone importance): 1.0
\item $\beta$ (heuristic importance): 2.0
\item Evaporation rate ($\rho$): 0.1
\end{itemize}

Results after 100 iterations:
\begin{itemize}
\item Best solution fitness: 0.247 (lower is better)
\item Convergence achieved at iteration 73
\item Component 1: Dual sourcing (65\% Supplier B, 35\% Supplier C)
\item Component 2: Single sourcing (100\% Supplier A)
\item Total annual cost: \$8,890
\item Risk score improvement: 22\% vs. single sourcing baseline
\item Computation time: 2.3 seconds
\end{itemize}

The ACO algorithm discovered that the slight increase in cost (\$130 vs. heuristic solution) was justified by the 7\% additional risk reduction through better supplier allocation patterns.

\subsection{Tier 4: The AI/ML/RL Augmentation Method (Deep Reinforcement Learning)}

A Deep Q-Network (DQN) approach treats sourcing decisions as a sequential decision-making problem where an agent learns optimal policies through interaction with a simulated supply chain environment.

\textbf{Reinforcement Learning Formulation}

\textbf{State Space:} The state $s_t$ encompasses:
\begin{align}
s_t = \{D_t, I_t, R_t, C_t, H_t\}
\end{align}
where:
\begin{itemize}
\item $D_t$: Current demand forecast vector
\item $I_t$: Inventory levels for all components
\item $R_t$: Supplier reliability scores (dynamic)
\item $C_t$: Current cost matrix
\item $H_t$: Historical disruption patterns
\end{itemize}

\textbf{Action Space:} For each component $i$, the agent selects:
\begin{align}
a_t^i = \{strategy_i, allocation_i\}
\end{align}
where $strategy_i \in \{single, dual\}$ and $allocation_i$ defines supplier percentages.

\textbf{Reward Function:} Multi-objective reward balancing cost efficiency and risk mitigation:
\begin{align}
R_t = -\lambda_1 \cdot Cost_t - \lambda_2 \cdot Risk_t + \lambda_3 \cdot Service\_Level_t \\
\end{align}

\lstinputlisting[language=Python, caption=Deep Q-Network for Sourcing Decisions]{.././codes-python/line89.py3.py}

\begin{figure}[htbp]
\centering
\includegraphics[width=\textwidth]{../figures-png/line89.fig4.png}
\caption{Deep Q-Network Architecture for Sourcing Decision Learning}
\end{figure}

\textbf{Concrete Example}

Training the DQN agent on a 52-week simulation with 3 components and 4 suppliers:

Training Configuration:
\begin{itemize}
\item Episodes: 1000
\item Learning rate: 0.001
\item Replay buffer size: 10,000
\item Target network update frequency: 100 episodes
\item $\epsilon$-decay: 0.995
\end{itemize}

Training Results:
\begin{itemize}
\item Final average reward: -2.34 (vs. -4.12 initial)
\item Convergence achieved at episode 743
\item Service level maintained: 98.7\%
\item Average weekly cost reduction: 12\% vs. static heuristic
\item Disruption recovery time: 1.8 weeks average
\end{itemize}

Learned Policy Insights:
\begin{itemize}
\item Critical components: 89\% dual sourcing frequency
\item Standard components: 34\% dual sourcing frequency  
\item Dynamic allocation adjustments during disruption periods
\item Predictive inventory building before peak seasons
\end{itemize}

The DQN agent learned to dynamically adjust sourcing strategies based on seasonal patterns, supplier reliability trends, and cost fluctuations, achieving superior performance compared to static optimization approaches.

\subsection{Tier 5: The Integrated Digital Twin (System-of-Systems Simulation)}

The Digital Twin paradigm transforms sourcing from static decision-making to dynamic, real-time optimization within a virtualized ecosystem that mirrors the entire supply chain network.

\textbf{System Architecture}

The Integrated Digital Twin creates a parallel virtual representation of the physical supply chain, enabling continuous ``what-if'' analysis and real-time decision optimization. This system-of-systems approach integrates:

\begin{itemize}
\item \textbf{Supplier Digital Twins:} Virtual models of each supplier's capacity, quality metrics, and operational status
\item \textbf{Demand Sensing Networks:} Real-time market intelligence and demand prediction models
\item \textbf{Logistics Network Twin:} Transportation, warehousing, and distribution center simulations
\item \textbf{Risk Monitoring Systems:} Geopolitical, weather, and economic risk assessment modules
\end{itemize}

\textbf{Interoperability Framework}

The system employs standardized APIs and data exchange protocols to ensure seamless information flow between subsystems:

\begin{itemize}
\item \textbf{Data Ingestion Layer:} IoT sensors, ERP systems, market data feeds
\item \textbf{Semantic Data Layer:} Ontology-based data harmonization across diverse sources
\item \textbf{Analytics Engine:} Machine learning models for pattern recognition and prediction
\item \textbf{Decision Support Interface:} Real-time dashboards and automated alerts
\end{itemize}

\begin{figure}[htbp]
\centering
\includegraphics[width=\textwidth]{../figures-png/line89.fig5.png}
\caption{System-of-Systems Digital Twin Architecture}
\end{figure}

\textbf{Simulation-Based What-If Analysis}

The Digital Twin enables comprehensive scenario planning through parallel universe simulations:

\textbf{Scenario 1: Supplier Disruption Impact}
\begin{itemize}
\item Simulate complete failure of primary supplier for Component A
\item Evaluate secondary supplier capacity and ramp-up time
\item Assess inventory depletion rates and stockout probabilities
\item Calculate financial impact and recovery timeline
\end{itemize}

\textbf{Scenario 2: Demand Surge Response}
\begin{itemize}
\item Model 300\% demand increase for emergency medical equipment
\item Analyze supplier scalability across the network
\item Optimize allocation between single and dual sourcing strategies
\item Evaluate transportation and warehousing bottlenecks
\end{itemize}

\textbf{Scenario 3: Cost Optimization Under Uncertainty}
\begin{itemize}
\item Simulate commodity price volatility scenarios
\item Test dynamic sourcing strategy switching thresholds
\item Evaluate hedging strategies and contract negotiations
\item Assess trade-offs between cost savings and supply security
\end{itemize}

\textbf{Concrete Example: Real-Time Sourcing Optimization}

MediTech's Digital Twin processes the following real-time inputs on March 15th, 2025:

\textbf{Input Data Stream:}
\begin{itemize}
\item Supplier A: 92\% capacity utilization, quality score 0.94
\item Supplier B: 78\% capacity utilization, quality score 0.89
\item Supplier C: Experiencing 2-day shipping delays due to weather
\item Market demand forecast: 15\% increase for next 6 weeks
\item Raw material costs: 8\% increase in semiconductor pricing
\end{itemize}

\textbf{Digital Twin Processing:}
\begin{enumerate}
\item \textbf{Real-time State Update:} All supplier models update capacity and performance metrics
\item \textbf{Demand Propagation:} Forecast changes cascade through all component requirements
\item \textbf{Constraint Analysis:} System identifies potential bottlenecks in 3 weeks
\item \textbf{Optimization Run:} Multi-objective solver recommends sourcing adjustments
\item \textbf{Impact Simulation:} 1000 Monte Carlo runs validate robustness
\end{enumerate}

\textbf{System Recommendations:}
\begin{itemize}
\item Shift Component 1 from 100\% Supplier A to 70\% A, 30\% B
\item Increase safety stock for Component 2 by 20\%
\item Activate pre-negotiated capacity expansion with Supplier B
\item Expected cost increase: 3.2\% vs. 12\% potential stockout costs
\item Implementation timeline: Adjust orders within 4 hours
\end{itemize}

\textbf{Validation Results:}
Over the following 6 weeks, the Digital Twin's recommendations resulted in:
\begin{itemize}
\item 99.7\% service level maintenance (vs. 94.3\% with static strategy)
\item Total cost increase: 3.4\% (within 0.2\% of prediction)
\item Zero stockout events (vs. 3 predicted stockouts with no action)
\item Supplier relationship strengthening through proactive communication
\end{itemize}

\subsection{Tier 7: The Human-AI Symbiotic Partnership (The Centaur Model)}

The Centaur Model transcends traditional automation by creating a synergistic partnership where human intuition and strategic thinking combine with AI's computational power and pattern recognition capabilities.

\textbf{Augmented Intelligence Framework}

Rather than replacing human decision-makers, the system amplifies human capabilities through:

\begin{itemize}
\item \textbf{Cognitive Augmentation:} AI provides real-time analysis while humans contribute strategic context and ethical considerations
\item \textbf{Explainable AI Integration:} All AI recommendations include transparent reasoning chains that humans can understand and modify
\item \textbf{Collaborative Learning:} The system learns from human feedback and domain expertise while teaching humans about hidden patterns
\item \textbf{Contextual Intelligence:} Humans provide nuanced understanding of relationships, politics, and intangible factors that AI cannot quantify
\end{itemize}

\textbf{Symbiotic Decision-Making Process}

The partnership operates through structured collaboration phases:

\textbf{Phase 1: AI-Driven Analysis}
\begin{itemize}
\item Process quantitative data at scale
\item Identify patterns and correlations
\item Generate multiple optimization scenarios
\item Highlight risk factors and trade-offs
\end{itemize}

\textbf{Phase 2: Human Interpretation}
\begin{itemize}
\item Apply strategic business context
\item Consider stakeholder relationships
\item Evaluate ethical implications
\item Incorporate market intelligence and intuition
\end{itemize}

\textbf{Phase 3: Collaborative Refinement}
\begin{itemize}
\item Humans modify AI recommendations based on qualitative factors
\item AI recalculates scenarios with human constraints
\item Interactive optimization through dialogue
\item Consensus building on final strategy
\end{itemize}

\begin{figure}[htbp]
\centering
\includegraphics[width=\textwidth]{../figures-png/line89.fig6.png}
\caption{Human-AI Symbiotic Partnership Architecture}
\end{figure}

\textbf{Explainable AI Implementation}

The system ensures transparency through multiple explanation mechanisms:

\textbf{Decision Tree Visualization:} AI shows the logical path from data inputs to sourcing recommendations using interactive decision trees.

\textbf{Feature Importance Ranking:} Quantitative breakdown of which factors most influenced each recommendation (cost 35\%, reliability 28\%, capacity 22\%, lead time 15\%).

\textbf{Counterfactual Explanations:} ``If supplier reliability had been 0.95 instead of 0.85, the recommendation would change to single sourcing.''

\textbf{Confidence Intervals:} All predictions include uncertainty ranges: ``Cost savings: \$45,000 ± \$8,000 (95\% confidence).''

\textbf{Concrete Example: Strategic Sourcing Crisis Response}

\textbf{Situation:} March 2025 - A major semiconductor shortage threatens MediTech's production schedule for Q2.

\textbf{AI Analysis (completed in 2 minutes):}
\begin{itemize}
\item Identified 47 alternative suppliers across 12 countries
\item Calculated cost implications: +18\% to +45\% price increases
\item Predicted delivery timelines: 3-8 weeks variability
\item Risk assessment: 73\% probability of continued shortages
\end{itemize}

\textbf{Human Strategic Input:}
Sarah Chen (Chief Procurement Officer) adds critical context:
\begin{itemize}
\item ``Supplier X in Taiwan has unused capacity but requires relationship building''
\item ``Regulatory approval for alternative components takes 6 weeks minimum''
\item ``Customer contracts include force majeure clauses we can invoke''
\item ``Board will accept 20\% cost increase to maintain delivery commitments''
\end{itemize}

\textbf{Collaborative Refinement:}
\begin{enumerate}
\item AI recalculates scenarios with human constraints
\item Sarah identifies Supplier Y (not in AI's database) through personal network
\item AI rapidly assesses Supplier Y's capacity and integration timeline
\item Joint decision: Split between 3 suppliers instead of AI's recommended 2
\item Human intuition: ``Something feels wrong about Supplier Z's pricing''
\item AI validates: Supplier Z shows statistical anomalies in pricing patterns
\end{enumerate}

\textbf{Final Strategy (Human-AI Collaborative):}
\begin{itemize}
\item Primary supplier (60\%): Existing Partner A with expedited delivery
\item Secondary supplier (25\%): Supplier Y through Sarah's network connection
\item Tertiary supplier (15\%): Qualified backup from AI's analysis
\item Regulatory strategy: Parallel approval process for 2 component alternatives
\item Customer communication: Proactive transparency about potential delays
\end{itemize}

\textbf{Outcome Validation:}
\begin{itemize}
\item Strategy implemented within 48 hours (vs. 2 weeks for traditional procurement)
\item Cost increase limited to 16\% (within board-approved threshold)
\item Zero production delays in Q2
\item Strengthened relationships with 2 new strategic suppliers
\item AI model updated with human domain knowledge for future scenarios
\end{itemize}

\textbf{Continuous Learning Loop:}
The system captured and integrated:
\begin{itemize}
\item Sarah's supplier network knowledge into relationship mapping
\item Regulatory timeline patterns for future planning
\item Customer communication strategies for crisis management
\item Anomaly detection improvements for supplier risk assessment
\end{itemize}

This symbiotic approach achieved superior outcomes by combining AI's computational power with human strategic thinking, relationship intelligence, and ethical reasoning, demonstrating the exponential value of human-machine collaboration over either approach alone.

\section{Practice Questions}

\textbf{Tier 1 Question:}
FormulaTech Manufacturing sources a critical component for their pharmaceutical equipment from two potential suppliers. Supplier Alpha offers the component at \$25 per unit with 95\% reliability, while Supplier Beta offers it at \$30 per unit with 88\% reliability. Annual demand is 10,000 units, and dual sourcing incurs an additional \$15,000 in management overhead. Formulate this as a multi-objective optimization problem with objectives of cost minimization and risk minimization. Define your decision variables, constraints, and solve for the optimal sourcing strategy assuming risk is quantified as (1-reliability) × annual demand value.

\textbf{Tier 2 Question:}
Implement the priority-based sourcing algorithm for a scenario with 4 components and 3 suppliers. Component criticality scores are [0.9, 0.7, 0.8, 0.6], supplier costs are [\$10, \$12, \$11] per unit, and reliability scores are [0.85, 0.92, 0.88]. Annual demands are [500, 800, 300, 1200] units respectively. Use the composite scoring function with weights (0.4, 0.4, 0.2) for cost, reliability, and capacity factors. Determine the sourcing strategy for each component and calculate the total annual cost.

\textbf{Tier 3 Question:}
Design an Ant Colony Optimization approach for a sourcing problem with 3 components and 4 suppliers. Set ACO parameters: $\alpha$ = 1.0, $\beta$ = 2.0, $\rho$ = 0.1, with 20 ants and 50 iterations. Component demands are [200, 350, 180] units, supplier costs range from \$8-\$15 per unit, and reliability scores are [0.82, 0.91, 0.87, 0.89]. Describe the pheromone update mechanism and calculate the expected convergence properties. What solution quality improvement would you expect compared to the greedy heuristic?

\textbf{Tier 4 Question:}
Develop a Deep Q-Network formulation for dynamic sourcing decisions over a 26-week horizon. Define the state space to include demand forecasts, inventory levels, and supplier performance metrics. Design the reward function to balance cost efficiency (\$50,000 average weekly spend), service level targets (98\% minimum), and supply chain resilience. Specify the neural network architecture and training parameters. How would you evaluate the learned policy's performance against traditional optimization approaches?

\textbf{Tier 5 Question:}
Design a Digital Twin system for real-time sourcing optimization across a network of 8 suppliers and 12 components. Specify the required IoT sensors, data integration APIs, and simulation engines needed to enable ``what-if'' scenario analysis. The system must handle 500 transactions per minute and provide optimization recommendations within 30 seconds. Calculate the computational requirements and evaluate three different disruption scenarios: single supplier failure, 25\% demand surge, and 15\% cost volatility across all suppliers.

\textbf{Tier 7 Question:}
Create a Human-AI symbiotic decision framework for strategic sourcing during supply chain crises. Design the explainable AI components that provide transparency for procurement managers making decisions under uncertainty. The system should handle qualitative factors like supplier relationships, regulatory constraints, and stakeholder politics alongside quantitative optimization. Describe how the system would learn from human expertise and adapt its recommendations. Evaluate the framework's effectiveness for a scenario involving geopolitical supply disruptions affecting 40\% of the supplier base.

\end{document}
